\documentclass{article}
\usepackage{amsmath,amssymb}
\title{Native dialgebra diagrams and varying-interface bodies}
\date{15 November 2026}
\begin{document}\maketitle
\paragraph{Status.} General presentation lemma, independently reviewed in the stated scope. Its purpose is to test actual transport of changing loop interfaces and directed native process squares, rather than add a diagram wrapper whose states have been forgotten. No feedback-preservation or existence theorem follows without the further hypotheses explicitly stated below.

\paragraph{Independent native input.} Let $P,Q:\mathcal S\to\mathcal K$ be pseudofunctors of bicategories, with the native dialgebra bicategory $\mathcal D$ already constructed. Its objects are $(X,f:PX\to QX)$; a native arrow is
\[
 (h,\alpha):(X,f)\to(Y,g),\qquad
 \alpha:Qh\,f\Rightarrow g\,Ph.
\]
A native two-cell $\sigma:(h,\alpha)\Rightarrow(k,\beta)$ is a state two-cell satisfying
\[
 \beta\,(Q\sigma\,f)=(g\,P\sigma)\,\alpha.
\tag{N}
\]
All equations include the specified pseudofunctor frames. No faithfulness of $P$ or $Q$ is assumed, and (N) is an actual process-cell equation. Native composition is the two-square paste. The state projection retains the full state two-cell, so it reflects equality of native two-cells and invertibility of an already admitted native two-cell. It does not create missing native squares.

\paragraph{Objects of the diagram presentation.} For an ordinary small shape $J$, the presentation of $\operatorname{Ps}(J,\mathcal D)$ consists of a state pseudofunctor $X:J\to\mathcal S$ and a directed lax transformation
\[
 f:PX\Rightarrow QX,\qquad
 f_a:Q(Xa)f_j\Rightarrow f_kP(Xa)
 \quad(a:j\to k).
\tag{1}
\]
Even though the diagram into $\mathcal D$ is a pseudofunctor, the cells $f_a$ need not be invertible. The pseudofunctor composition comparison lives one dimension higher: it is the invertible state compositor $Xb\,Xa\Rightarrow X(ba)$ satisfying (N) for the pasted transition square and $f_{ba}$. That equation is precisely the lax-transformation composition equation for (1). Units give its unit equation. Thus no directed transition square is inadvertently made invertible.

Conversely, given $X$ and (1), define the diagram at $j$ by $(Xj,f_j)$, at $a$ by $(Xa,f_a)$, and use the state pseudofunctor's constraints as the native diagram constraints. The lax equations prove their native admission; their inverses are admitted by rearranging (N) with the invertible state comparison. The diagram pentagon and units then follow from those of $X$, since the native projection detects equality of these admitted constraints. This proves the object correspondence constructively.

\paragraph{Full arrows and witnesses.} A pseudonatural diagram arrow corresponds to a pseudonatural state transformation $h:X\Rightarrow Y$ and a modification
\[
 \alpha:(Qh)f\Rightarrow g(Ph).
\tag{2}
\]
Its component is the actual native square $\alpha_j:Qh_j f_j\Rightarrow g_jPh_j$. Along $a:j\to k$, the state naturality cell is
$\sigma_a:Ya\,h_j\Rightarrow h_k\,Xa$. The two composite native squares are, schematically,
\[
 QYa\,Qh_j f_j\Rightarrow QYa\,g_jPh_j
 \Rightarrow g_kPYa\,Ph_j
\]
and
\[
 Qh_k\,QXa f_j\Rightarrow Qh_k f_kPXa
 \Rightarrow g_kPh_k\,PXa.
\]
Condition (N) for $\sigma_a$ says that the second paste precomposed by $Q\sigma_a f_j$ equals the first paste postcomposed by $g_kP\sigma_a$. This is exactly the modification equation for (2). Hence it accounts for the full interaction of state naturality and transition cells, not merely pointwise components. The pseudonaturality axioms are the state transformation axioms with these admitted equations.

A modification between two such diagram arrows is a state modification $\tau:h\Rightarrow h'$ whose components satisfy (N) with $\alpha_j,\alpha'_j$. Its shape-arrow equation is the state modification equation. These are precisely the original native witnesses and their actual compatibility conditions. Composition is obtained by pasting (2); its components are the original two-square composition. Associators and units agree under the extraction/insertion described above. With the chosen frames this gives the resulting biequivalence of bicategories, not a quotient by state reparameterization.

\paragraph{Lax-body enhancement.} The same calculation allows the diagram arrow to be directed lax rather than pseudonatural. Then $h$ is a directed lax state transformation and its cells $\sigma_a$ need not be invertible. The displayed equation for the two composite squares uses $\sigma_a$ only forward, so no inverse is required. Conversely (2) admits each such $\sigma_a$ as a native diagram naturality cell. Its composition/unit axioms are again the lax state transformation axioms. This yields the corresponding presentation of the fixed-shape lax-body diagram bicategory used in synchronization coreflection. One must not confuse laxity of $f$ in (1) with laxity of the map $h$ between two such bodies; both dimensions are retained.

Restriction along $u:I\to J$ acts simultaneously on $X,f,h,\alpha$ and their witnesses. A shape transformation $\eta:u\Rightarrow v$ uses $X\eta$ together with the transition cells $f_{\eta_i}$. These give the actual native diagram arrow, rather than a bare state comparison. Thus total reverse-shape presentation retains the changing-interface data as well. Naturality and composition follow from the same state pseudofunctor and lax-body equations.

\paragraph{Application to state feedback.} In the original relative state construction, $P=P_A$ and $Q=Q_B$ send a retained state $X$ to its actual pure action on external interfaces $A,B$. Formula (1) therefore presents a family with changing loop interfaces $P_AX_j,Q_BX_j$ and native transition squares along each shape arrow. It is not a family whose outer state is silently fixed, and it is not a list of marginal transition laws. The original pure embedding and its action comparison frames remain part of $P_A,Q_B$.

If selected native lax bilimits exist in $\mathcal D$, their synchronization evaluates precisely this retained diagram. Cone creation still requires its independent state-limit and $Q$-preservation hypotheses; the presentation lemma supplies no such limits. If a feedback evaluator is given as a genuine pseudofunctor on the entire native dialgebra, postcomposition acts on these full diagrams. Strong commutation with synchronization requires actual preservation of the selected cones; a directed comparison requires its appropriate cone transport. Objectwise feedback functions or independent naturality at fixed states do not provide (2) or (N) and cannot be substituted.

For the canonical finite tower, join-then-synchronize remains an already proved route: its joining pseudofunctor first transports these full native diagrams to the level where cone creation is available. The present lemma explains the varying-interface input exactly; it does not strengthen that result into preservation of unavailable higher-level limits or an instantaneous trace.

\paragraph{Relation to established machinery.} This is a direct unpacking of pseudodiagrams in a native lax dialgebra, using the established bicategories of lax transformations and modifications. It is not claimed as a new general theory of lax comma objects. The project-specific check is the direction (N), the difference between diagram compositors and transition squares, and retention of the original pure-action interfaces without a faithful extension assumption. These checks have received independent written-proof review.
\end{document}
